\begin{table}[htbp]
\centering
\footnotesize
\setlength{\tabcolsep}{4pt}
\caption{更新型比值样本的二项稀释与州固定效应拟合。表 A 列出每个保留概率 $q$ 下原始尺度 $C_t$ 的斜率、隐含稀释斜率 $c_q=[b(q)-b(1)]/(1/q-1)$、截距变化 $a(q)-a(1)$ 以及 $C_t\ge50$ 子集 ($n=828$) 的 $c_q$。表 B 列出各模型下 $1/C_t$ 的斜率与平衡比。括号内为 95\% 双向自助法置信区间。}
\label{tab:s11}
\begin{tabular}{lllll}
\toprule
\multicolumn{5}{l}{\textit{A. 受控二项稀释实验 ($n=1748$)}}\\
$q$ & $b(q)$ & $c_q$ & $\Delta a$ ($\times10^{-3}$) & $c_q$, $C_t\ge50$ \\
\midrule
0.8 & 0.676 [0.466, 0.903] & 0.463 [0.061, 0.909] & $-$0.32 [$-$1.74, 1.18] & 0.472 [$-$0.521, 1.374] \\
0.6 & 0.880 [0.577, 1.212] & 0.479 [0.157, 0.875] & $-$0.87 [$-$3.65, 1.81] & 0.456 [$-$0.167, 1.082] \\
0.5 & 1.062 [0.724, 1.447] & 0.502 [0.242, 0.850] & $-$1.36 [$-$4.32, 1.60] & 0.456 [$-$0.022, 0.909] \\
0.4 & 1.321 [0.959, 1.865] & 0.507 [0.254, 0.909] & $-$1.85 [$-$6.36, 2.04] & 0.458 [0.059, 0.908] \\
0.3 & 1.745 [1.275, 2.384] & 0.508 [0.315, 0.764] & $-$2.52 [$-$7.53, 1.35] & 0.470 [0.081, 0.856] \\
0.2 & 2.439 [1.894, 3.152] & 0.470 [0.339, 0.646] & $-$2.33 [$-$8.70, 2.81] & 0.474 [0.167, 0.731] \\
0.1 & 4.082 [2.858, 5.481] & 0.391 [0.260, 0.552] & 0.11 [$-$9.02, 9.98] & 0.487 [0.212, 0.779] \\
\midrule
合并稀释斜率 $c$ & -- & 0.413 [0.305, 0.544] & -- & 0.483 [0.253, 0.717] \\
泊松基准系数 $b_P$ & -- & 0.442 & -- & 0.426 \\
\bottomrule
\end{tabular}

\bigskip
\begin{tabular}{lll}
\toprule
\multicolumn{3}{l}{\textit{B. 州固定效应模型 ($n=1748$, 51 个辖区)}}\\
模型设定 & 斜率 $b$ & 平衡比 $b/a$ \\
\midrule
合并横截面回归 (Pooled) & 0.560 [0.401, 0.760] & 37.4 [21.6, 67.9] \\
州固定效应，算术均值 $a_s$ & 0.379 [0.117, 0.663] & 17.3 [6.0, 77.1] \\
州固定效应，中位数 $a_s$ & -- & 21.7 [8.0, 157.6] \\
州固定效应，行加权均值 $a_s$ & -- & 18.4 [6.6, 79.1] \\
州固定效应，病例加权均值 $a_s$ & -- & 24.6 [8.4, 87.4] \\
州固定效应，分母规模 $D_t$ & 0.379 [0.117, 0.663] & 17.3 [6.0, 77.1] \\
仅周固定效应 (Week intercepts) & 0.508 [0.220, 0.963] & -- \\
州与周双向固定效应 & 1.02 [$-$1.05, 3.06] & -- \\
横截面与州内均值配对比值 & -- & 2.16 [0.59, 5.15] \\
\bottomrule
\end{tabular}
\end{table}
